
\section{The Peira framework}\label{sec:background}

Peira is a framework for systematic conduct of experimental studies for the purpose of evaluating the appropriateness of conceptual modeling language vocabularies \cite{LZMK.24}. An instance of a Peira experiment is meant to evaluate a set $V$ of \textit{signifiers} representing the vocabulary of the language to be evaluated. To do so, a sample of descriptions $E$ is taken from one or more application domains where the language is meant to be used. Further, a set of $\mathbf{D} = D \cup R$ of distinguished elements $D$ from each domain, combined with a set of $R$ n-tuples thereof, are defined as a sample of the Universe of Discourse (UoD) relating to the domain and deemed as worthy of being modeled using the vocabulary. Domain elements from $\mathbf{D}$ under specific descriptions are referred to as \textit{subjects} from a set $S = \mathbf{D} \times E$. A set $P$ of \textit{raters} are also identified as a sample (or a proxy thereof) from a population of potential users of the language whose vocabulary is in question. The role of the raters is to study each description  $e\in E$ and classify each subject $s\in S$ associated with $e$ to none, one, or more of the terms in vocabulary $V$.

As an example, consider a simple domain-specific language for describing enrollment and attendance of students to courses and lectures, using concepts \textit{course}, \textit{student}, \textit{lecture}, (entities / unary concepts) \textit{enrolls}, \textit{includes lecture}, and \textit{attends lecture} (relationships / binary concepts). Our evaluation question is if the vocabulary $V = $ \{``course'', ``student'', ... \} --- the English words representing the above concepts --- is appropriate for describing situations involving students enrolling courses, attending lectures, etc. A set of descriptions $E$ is prepared of real or contrived cases of students taking courses, attending lectures etc. Consider that in one of them $e\in E$, Alice decided to take Statistics only to find that it is hard to follow and that she fell asleep during the lecture on the Central Limit Theorem (CLT). Within the context of the case, elements of  $D$ include ``Alice'', ``Statistics'', ``Lecture on CLT'' and elements of $R$ include tuples such as ``(Alice, Statistics)'', ``(Alice, Lecture on CTL)'', and ``(Statistics, Lecture on CTL)''. These elements combined constitute subjects $s\in S$ under the description of Alice's case $e$. These subjects are to be classified under none, one, or more of the elements in $V$ by a sample of participants who are likely to use the vocabulary, such as the maintainers and analysts/instantiators of a software product line for course enrollment and attendance tracking.

The result of the classification exercises is a dataset of triples $(p,s,v)$ --- rater $p$ classified subject $s$ under term $v$ --- which Peira proposes can be the basis for several metrics of assessing different qualities of the vocabulary, assuming representativeness of descriptions, subjects, and raters. For example, when a term $v$ is rarely or never featured in any triple, this can be construed as evidence of construct excess: $v$ could be dropped from the language. Reversely, subjects $s\in S$ not appearing in any triples could be evidence of construct deficit: while the experimenters thought that $s$ is worthy of being modeled, hence included in $S$, the participants could not classify it under any of the available terms. The authors offer concrete formulae for quantitatively expressing those issues \cite{LZMK.24}. 

For simplicity, we will, here, focus on Peira metrics that compare participant responses with the authoritative classifications, i.e., the complete set of $(p_a,s,v)$ triples that the language designers believe are valid classifications. Following \cite{LZMK.24} let function $n: P\times S \times V \mapsto \{0,1\}$, be $n(p,s,v) = 1$ if participant $p$ has offered classification $(p,s,v)$, and $n(p,s,v) = 0$ otherwise. Let now $p_a$ be a participant from $P$ that holds the authoritative response. Define the following for $p\neq p_a$:
\[ 
\mathit{acc}(p,s,v)= \left\{
\begin{array}{ll}
	\mathbf{1}, \:\mathrm{if}\,n(p_a,s,v) = 1\,\mathrm{and}\,n(p,s,v) = 1\\
	\textbf{0}, \:\mathrm{otherwise} \\
\end{array} 
\right. 
\]
\[
\mathit{def}(p,s,v)= \left\{
\begin{array}{ll}
	\mathbf{1}, \:\mathrm{if}\,n(p_a,s,v) = 1\,\mathrm{and}\,n(p,s,v) = 0\\
	\textbf{0}, \:\mathrm{otherwise} \\
\end{array} 
\right. 
\]
\[
\mathit{exc}(p,s,v)= \left\{
\begin{array}{ll}
	\mathbf{1}, \:\mathrm{if}\,n(p_a,s,v) = 0\,\mathrm{and}\,n(p,s,v) = 1\\
	\textbf{0}, \:\mathrm{otherwise} \\
\end{array} 
\right. 
\]

We further use $\mathit{acc}(\cdot,\cdot,v)$, $\mathit{def}(\cdot,\cdot,v)$, and 
$\mathit{exc}(\cdot,\cdot,v)$ to represent the marginal totals for each term over all subjects and all participants, e.g., $\mathit{acc}(\cdot,\cdot,v) = \sum_{p\in {P \setminus {p_a}}}\sum_{s\in S} \mathit{acc(p,s,v)}$. According to Peira, the quantities offer raw measures of the {\em accuracy}, {\em scope deficiency} and {\em scope excess} of a given term. Specifically, for a specific term $v$, $\mathit{acc}(\cdot,\cdot,v)$ represents the number of ratings in which participants and designers are in agreement with respect to term $v$, $\mathit{def}(\cdot,\cdot,v)$ the number of authoritative ratings that were not picked by participants (``false negatives''), and $\mathit{exc}(\cdot,\cdot,v)$ the participant ratings that are in excess of the authoritative ones (``false positives''). It follows that, in the more familiar terms of precision and recall:

\noindent $\mathrm{prec}^{pc}_V(v) = acc(\cdot,\cdot,v)/(acc(\cdot,\cdot,v)+exc(\cdot,\cdot,v))$

\noindent $\mathrm{rec}^{pc}_V(v) = acc(\cdot,\cdot,v)/(acc(\cdot,\cdot,v)+def(\cdot,\cdot,v))$

Seeing participants as classifiers of labeled data (the authoritative classifications), the above describe the \textit{precision-per-concept} and \textit{recall-per-concept} of the corresponding classification exercise. Concepts for which these metrics are comparatively low are subjects for investigation by designers.

Likewise, $\mathit{acc}(p,\cdot,\cdot)$, $\mathit{def}(p,\cdot,\cdot)$, and $\mathit{exc}(p,\cdot,\cdot)$ offer corresponding measures characterizing the input of a specific participant compared to the authoritative. These ratings measure the degree to which the participant is in general agreement with the designers vis-a-vis the vocabulary, and characterize the performance of that participant. In terms of precision and recall again:

\noindent $\mathrm{prec}^{pp}_V(p) = acc(p,\cdot,\cdot)/(acc(p,\cdot,\cdot)+exc(p,\cdot,\cdot))$

\noindent $\mathrm{rec}^{pp}_V(p) = acc(p,\cdot,\cdot)/(acc(p,\cdot,\cdot)+def(p,\cdot,\cdot))$

These metrics are called \textit{precision-per-participant} and \textit{recall-per-participant} and are used in Peira to also evaluate the language in a way that is amenable to statistical inferences. Below we will use these metrics to compare participant performance.

While Peira is agnostic to how experimenters acquire ratings in practice, in their experiments, the authors have utilized questionnaire-style instruments in which participants, after being shown training videos, are presented a case, followed by a presentation of subjects below each of which a multi-choice inventory of vocabulary terms (plus the ``None'' option) are offered for participants to classify the subject under one or more terms \cite{LZMK.24}. In the experimental study we followed the same rating acquisition style, as we describe below.